% ===================================================================== Chapter 10
\chapter{Thermal Field Transfer}\label{ch:mapping}

\section{CFD-to-Mechanical Workflow}
The converged medium-mesh solution was transferred to ANSYS Mechanical without re-solving any thermal problem: the Fluent solid temperature
field is imported directly into Static Structural as a body temperature (decision D-039). The source case and data files were hash-verified
before and after every structural run and are unchanged. Nineteen read-back checks confirmed that the source is the converged second-order
solution with the correct solid material and a temperature range inside the property tables.

\section{One-Way Coupling}
The coupling is one-way: CFD temperature to structure, with no feedback of deformation to the flow. The analytical thermal growth of the bore radius,
34.94~\textmu m, changes the flow area by only 0.70\,\%, so the flow field is not affected by the deformation. A two-way or fluid--structure interaction analysis would not change the conclusions and was not performed.

\section{Temperature Mapping}
The first method tried, an External Data point cloud of cell and face centroids with triangulation, failed validation: it flagged 77,088
of the 111,216~nodes of the structural mesh then in use as outside the cloud, left 27,696~nodes without a temperature and produced errors of up to 4.0~K. Kriging and
radial-basis interpolation failed in batch mode, and a Workbench Fluent-system route held only the initial 300~K field. These attempts are
kept as evidence. The accepted method is mesh-based External Data:
\begin{itemize}
\item the master mesh is the Fluent solid mesh itself, rebuilt from the case file as an MAPDL archive;
\item the data are Fluent's own node temperatures (48,048~nodes), exported in EnSight format and keyed by node number;
\item Mechanical maps them with the Bucket Volume search and Shape Function interpolation; the outside option is Nearest Node.
\end{itemize}
The 9,408 structural nodes on the outer surface that lie in the 0.038~mm chord gap between the true circle and the 48-sided CFD polygon
receive shape-function values of the adjacent source element (verified), not a nearest-node copy.

\section{Mapping Validation}
Table~\ref{tab:mapping} and Figures~\ref{fig:mapside}--\ref{fig:mapaxial} summarise the checks of the transfer; ``validation'' here means a numerical comparison with the Fluent source field and an
independent reference, not an experimental validation. All 108,252 structural nodes received a
temperature; the mapped range, 423.84--562.56~K, matches the source node range of 423.84--562.54~K. The mid-span through-wall temperature
difference (7.580~K mapped against 7.581~K in Fluent) and the volume mean (525.520~K against 525.483~K) are preserved. In Mechanical the imported field is displayed
in degrees Celsius, 150.69--289.41\,\textdegree C.

%% generated by gen_tables.py - RE-ANALYSIS 2026
\begin{table}[htbp]
\centering
\footnotesize
\caption{Numerical checks of the CFD-to-Mechanical temperature transfer (Section 7A)}
\label{tab:mapping}
\begin{tabular}{>{\raggedright\arraybackslash}p{0.42\linewidth}>{\raggedright\arraybackslash}p{0.48\linewidth}}
\toprule
Check & Result \\
\midrule
Source field & converged medium solution; Fluent node temperatures 423.84--562.54 K (48,048 nodes) \\
Coverage & 108,252 / 108,252 structural nodes mapped; 0 unmapped \\
Mapped range & 423.84--562.56 K (identical in LC1 and LC2) \\
Error (A) against exact shape-function evaluation & maximum 0.081 K inside the source mesh; mean 0.002 K \\
Error (B) against independent finite-volume reference & maximum 2.57 K at the bore within $\sim$11 mm of the inlet; $\le$ 0.25 K over $z$ = 14--586 mm; mean 0.039 K \\
Mid-span through-wall $\Delta T$ & mapped 7.580 K against Fluent 7.581 K \\
Volume-mean temperature & mapped 525.520 K against Fluent 525.483 K \\
Parametric cases (9B-2) & all 7 cases passed; error (A) 0.086--0.154 K \\
\bottomrule
\end{tabular}
\par\smallskip\parbox{\linewidth}{\footnotesize\raggedright Source: \texttt{07\_Thermal\_Analysis/THERMAL\_MAPPING\_NOTES.md}, \texttt{CFD\_TO\_MECHANICAL\_MAPPING\_AUDIT.md}, \texttt{PROJECT\_STATE.md} \S13.1 and \S19.2.}
\end{table}


\begin{figure}[htbp]
\centering
\includegraphics[width=0.68\linewidth]{F7A_03_side_by_side_and_difference_nt.png}
\caption{Fluent source temperature (top), temperature mapped onto the Mechanical nodes (middle) and difference from an independent
finite-volume reference (bottom), circumferential means [K].}
\label{fig:mapside}
\src{\provdata}
\end{figure}

\begin{figure}[htbp]
\centering
\begin{subfigure}[t]{0.56\linewidth}\centering\includegraphics[width=\linewidth,height=52mm,keepaspectratio]{F7A_04_axial_wall_temperature_nt.png}\caption{Inner and outer wall temperature and through-wall $\Delta T$ along the duct}\end{subfigure}\hfill
\begin{subfigure}[t]{0.42\linewidth}\centering\includegraphics[width=\linewidth,height=52mm,keepaspectratio]{F7A_05_through_wall_midspan_nt.png}\caption{Through-wall profile at mid-span}\end{subfigure}
\caption{Mapped temperatures compared with the Fluent source values.}
\label{fig:mapaxial}
\src{\provdata}
\end{figure}


\section{Mapping Error}
Two error measures are distinguished. Error (A) compares the mapped value with the exact shape-function evaluation of the source node
field: at most 0.081~K inside the source mesh, mean 0.002~K. This is the error of the transfer itself. Error (B) compares the mapped value
with an independent finite-volume reference built from the cell and face values: at most 2.57~K, at the bore within about 11~mm of the inlet,
and at most 0.25~K over $z$ = 14--586~mm. Error (B) reflects Fluent's reconstruction of node values from cell values in the region of the
steepest gradients; it is a property of the source data, identical on every structural mesh. Its effect was bounded at $\pm$9.8~MPa on the
LC1 inlet peak stress and 1.6\,\% on the LC2 peak; interior values are unaffected. Every parametric case was mapped with the same method and
passed the same checks, with error (A) between 0.086 and 0.154~K.

\section{Reference Temperature}
The stress-free reference temperature is $T_{ref}$ = 300~K, equal to the inlet air temperature (decision D-041).
The thermal load is the full spatial field $\Delta T(x,y,z) = T - T_{ref}$, which ranges from 123.8 to 262.6~K; neither the analytical mean
temperature nor the CFD maximum is used as a uniform substitute. The reference value was read back from every solver input file.

\section{Temperature-Dependent Material Properties}
The structural material uses $E(T)$ from the VDM table (204~GPa at 20\,\textdegree C to 180~GPa at 400\,\textdegree C) and the secant $\alpha(T)$
from 12.8 to $14.8\times10^{-6}$/K, re-referenced by the solver from the 70\,\textdegree F datum of the data to $T_{ref}$. Every node lies
inside the tabulated temperature range, so no property is extrapolated. $S_y(T)$ is interpolated at the node temperature only for the yield
assessment.
